8. Conclusion
This study examined probabilistic electricity demand forecasting under extreme cold-weather stress through a retrospective case study of the January 6–8, 2014 Polar Vortex in the PJM RTO system. Using PJM RTO hourly metered load data spanning 2010–2014, with a training period of 2010–2013 and a held-out test year of 2014, Linear Regression and GBoost point forecast models, a QR-GBT probabilistic model, three naive load-memory baselines, and an external PJM Day-Ahead benchmark were evaluated. The January 6–8, 2014 Polar Vortex produced a near-annual-peak RTO load of 140,510.2 MW at Jan 7 18:00 EPT, equal to 99.18% of the 2014 annual peak of 141,677.9 MW recorded on Jun 17 17:00 EPT. All weather-informed model results use same-hour ERA5 retrospective reanalysis weather input and reflect an idealized retrospective weather-information setting rather than an operational forecast-weather setting.
Among evaluated point forecast models, GBoost achieved the lowest errors across both the full-year 2014 test set (MAE 721 MW, RMSE 994 MW) and the January 6–8 vortex window (MAE 1,705 MW, RMSE 2,037 MW). All models exhibited degraded performance during the vortex window relative to their full-year rates, with the seasonal naive baselines showing the sharpest increases in error. Linear Regression, despite using the same ERA5 retrospective weather inputs as GBoost, produced a vortex-window MAE of 4,248 MW, compared with 1,705 MW for GBoost, suggesting that the linear specification was less able than GBoost to represent event-window behavior under extreme cold conditions in this setting. GBoost full-year MAE (721 MW) and vortex-window MAE (1,705 MW) were lower than the corresponding PJM Day-Ahead figures (1,937 MW and 3,148 MW respectively); however, these comparisons are retrospective in character, reflecting the informational advantage conferred by same-hour ERA5 reanalysis inputs, and should not be interpreted as evidence of operational superiority over PJM Day-Ahead.
The QR-GBT probabilistic model achieved a full-year q50 MAE of 761 MW and a nominal 90% prediction interval empirical coverage rate of 86.8%, and a nominal 98% coverage rate of 97.3%, across the 2014 test year. Disaggregation by evaluation window, however, reveals that these full-year aggregate metrics conceal substantial event-window undercoverage. During the January 6–8 vortex window, nominal 90% PI coverage fell to 66.7% and nominal 98% PI coverage fell to 84.7%. The top 1% of 2014 test hours by observed load showed 90% PI coverage of 63.2% and 98% PI coverage of 82.8%, consistent with the vortex-window findings. At the event peak, the observed load exceeded the post-rearrangement QR-GBT q99 estimate of 139,585 MW, leaving the event peak outside the nominal 98% prediction interval. A comparison of the January vortex window and the June 16–18 summer near-peak window — which contains the 2014 annual peak of 141,678 MW — indicates that similar peak-load magnitude does not imply similar probabilistic forecast performance: the winter window showed approximately twice the q50 MAE and mean pinball loss of the summer window, and substantially lower prediction interval coverage across both nominal levels. These results indicate that point forecast accuracy and probabilistic interval calibration are distinct properties, and that acceptable aggregate calibration does not guarantee adequate coverage during tail-risk conditions. Quantile crossing was detected in approximately 66% of full-year test hours; post-hoc monotonic rearrangement was applied as a correction, though this remains a methodological limitation rather than an architectural solution.
Taken together, the findings of this case study suggest that gradient boosted models with ERA5 retrospective weather inputs can substantially reduce point forecast error relative to naive baselines and linear models, including during a near-annual-peak cold-weather stress event, but that probabilistic coverage under the same conditions degrades in ways that full-year aggregate calibration metrics do not adequately reflect. This study does not establish operational deployment readiness or operational superiority over PJM Day-Ahead; both conclusions would require evaluation using operational forecast-weather inputs and a prospective validation design. Several directions for future work follow from these findings. First, pairing the modeling framework with ensemble NWP forecast archives or operational forecast-weather inputs would quantify the performance degradation attributable to weather forecast uncertainty and provide a more operationally grounded assessment. Second, extending evaluation across additional test years, grid systems, and extreme-event types — including heat waves and other high-impact weather regimes — would assess whether the patterns observed here generalize beyond the PJM 2014 case study. Third, directly characterizing training-distribution support for extreme cold events, using multi-year weather and load percentile profiles, would clarify the extent to which the observed vortex-window degradation is attributable to distribution shift. Fourth, probabilistic models with architecturally enforced joint quantile monotonicity, or alternative calibration frameworks, may reduce the reliance on post-hoc corrections and improve tail-risk coverage under stress conditions.